\section{模型的评价与推广}

\subsection{模型评价}

本文模型的主要优点在于任务侧与能源侧采用了统一时间口径。分钟级持续时间通过实际小时重叠量进入 GPU-hour 和功率计算，避免了将所有跨小时任务按完整小时重复占用。任务唯一性、实时即时开工、有向网络时延、截止时间、GPU 容量、IT 功率和设施功率被同时核验；能源侧则区分新能源直供、新能源充电、电网充电、外送和弃电，使成本、碳排放和新能源利用率能够从同一能源流中重新计算。

预测过程采用严格时间顺序切分，并对 Huber 候选设置零方差、矩阵秩、有限值和收敛门禁。固定起点递归预测避免了常见的时间泄漏问题，季节基线与增强模型之间的选择规则也在测试前确定。Huber 损失能够降低极端负荷对回归参数的影响\cite{huber1964robust}，而固定时间起点验证符合时间序列预测的基本评价原则\cite{hyndman2018forecasting}。

调度与储能计算采用确定性有界启发式，能够在大规模任务和共享墙钟预算下快速构造完整可行方案。每次局部移动均从现有占用中移除原任务，并在接受候选前复核容量；储能状态从唯一初始 SOC 前向递推，终端状态和分源守恒被逐时验证。该设计提高了结果的可复现性，且固定随机种子、阶段截止和候选顺序使同一输入下的执行路径保持确定。

模型的局限性主要来自搜索范围。理论滚动 $\varepsilon$-约束 MILP 并未运行，问题二和问题四只检查按 GPU-hour 排序的有界弹性任务子集，储能规则还采用了不低于初始 SOC 的保守状态子集。因此，正式结果只能说明候选在既定情景和约束下可行，不能代表附件完整可行域的全部优化潜力。固定文件中的候选比较表也不构成完整非支配前沿，相关多目标结论应被理解为有限候选之间的权衡\cite{miettinen1999multiobjective}。

此外，问题三累计放电量为 $0.0$ MWh，且有储能减无储能的成本变化为 $304223.92909026146$ 元，表明价格四分位规则在本次数据下未产生成本改善。该结果并非模型失败，而是说明“可行储能动作”与“经济性改善”不能被等同。服务质量指标也只由等待余量、即时开工、时延与截止约束等题内变量构造，其 $99.94443144071555\%$ 仅为代理值，不代表真实用户满意度。

\subsection{模型推广}

该框架可推广至云计算集群、边缘计算网络和多园区数据中心。若任务持续时间和 GPU 需求仍可被提前估计，只需替换任务类型功率映射、资源容量和有向时延矩阵，即可复用重叠计量、可行初始方案和局部改进机制。对于更大的任务规模，可采用滚动窗口保持跨窗口任务的固定后续占用，并继续使用完整候选证书防止局部提交破坏全时域可行性。

能源侧可进一步引入更完整的设备模型。若附件提供储能退化成本、自放电、循环寿命和动态效率，可将等效循环量与寿命成本纳入候选评分；若提供区域间线路容量、损耗和潮流数据，则可把当前独立区域平衡扩展为网络化电力流模型。此时仍应保留能源来源拆分，避免把电网购电转售错误计入新能源利用量。储能与可再生能源协同建模可参考综合能源系统的一般处理方式\cite{lindenberger2017energy}。

预测模块可在获得更长历史序列后扩展为概率预测，为每个区域—类型组合输出需求区间而非单点值。相应调度可采用鲁棒优化或情景优化，在不确定持续时间、新能源出力和价格条件下保留容量裕度。然而，新增模型必须继续遵守时间顺序验证，参数选择不能读取最终测试区间，概率覆盖率也必须由真实留出数据计算。

在工程部署中，模型还应连接滚动数据校验、异常告警和调度回退机制。若输入缺失、任务候选集合为空、储能边界无效或完整残差复核失败，应保留最近的完整可行方案并输出结构化失败状态。对于价格、碳强度或新能源出力变化，应重新运行任务与能源规则，而不能按比例改写既有结果。由灵敏度实验可知，当前目标对价格尺度较敏感，因此电价质量、售电结算口径和区域接入边界应被列为部署前的重点核验对象。
